\section{Finite trace tests and a global denominator budget}
\label{sec:finite}

The exact formula of the previous section is useful only when its trace
condition can be checked. A finite list of observed primes provides upper
bounds; an exact equality needs a certificate for the lower bound too.
The following finite test states explicitly when a finite list suffices.
It is a consequence of matrix algebra, and we make no priority claim for it.

\begin{theorem}[A finite certificate from group generators]
\label{thm:finite}
Let \(G=\langle S\rangle\), where \(S\) is a finite set of group generators.
Let \(\rho,\sigma:G\to\mathrm{GL}_d(O)\), where \(O\) is a discrete valuation
ring, \(d\geq1\), and \(\operatorname{span}_O\rho(G)=M_d(O)\).
Put \(m=d^2\). There are words \(g_1=1,g_2,\ldots,g_m\), each of length
at most \(m-1\) in \(S\), for which
\[
 B_i=\rho(g_i)\qquad(1\leq i\leq m)
\]
is an \(O\)-basis of \(M_d(O)\). For every \(r\geq1\), the two
representations are equivalent modulo \(\pi^r\) if and only if their
traces agree modulo \(\pi^r\) on the finite set
\[
 \mathcal T=\{g_i g_j:1\leq i,j\leq m\}
 \ \cup\
 \{s g_i g_j:s\in S,\ 1\leq i,j\leq m\}.
\]
All test words have length at most \(2d^2-1\).
If both integral matrix spans are full, their exact congruence depth is
\[
 a=\min_{w\in\mathcal T}v_\pi
       \bigl(\operatorname{tr}\rho(w)-\operatorname{tr}\sigma(w)\bigr).
\]
\end{theorem}

\begin{proof}
Over the residue field let \(V_j\) be the span of images of positive
words of length at most \(j\). Its dimension starts at one. If
\(V_{j+1}=V_j\), left multiplication by each generator preserves it.
These multiplications are invertible, so their restrictions to the
finite-dimensional space are surjective; inverses also preserve it.
Thus it contains every group image, so is already \(M_d(k)\).
It follows that \(V_{m-1}=M_d(k)\). Choose a basis of word images
including the identity and lift it to \(O\). Its coordinate determinant
is a unit, giving the asserted integral basis.

Work over \(R=O/\pi^r\). The trace pairing
\((X,Y)\mapsto\operatorname{tr}(XY)\) on \(M_d(R)\) is perfect.
The matrix \(\Gamma=(\operatorname{tr}(B_iB_j))_{i,j}\) has unit
determinant. The pair tests say that the same is true for the corresponding
matrix on \(B'_i=\sigma(g_i)\). Consequently the \(B'_i\) also form a
basis, and there is a unique linear isomorphism
\(\Phi:M_d(R)\to M_d(R)\) taking \(B_i\) to \(B'_i\).
It sends the identity to the identity.

Write \(\rho(s)B_i=\sum_k c_kB_k\).
Pair this equation with each \(B_j\), and use the triple tests. Perfectness
of the pairing on the \(B'_j\) implies
\[
 \sigma(s)B'_i=\sum_k c_kB'_k
              =\Phi(\rho(s)B_i).
\]
Thus \(\Phi\) intertwines left multiplication by all generators, and hence
their inverses and all group elements. Applying this to the identity
gives \(\Phi(\rho(g))=\sigma(g)\) for every \(g\).
As group images span, \(\Phi\) is multiplicative. The elementary
matrix-unit argument in the trace-rigidity proof makes it inner, proving
equivalence. The reverse implication is invariance of trace.
Taking the largest valid \(r\), with zero and infinity allowed, proves
the last assertion.
\end{proof}

\paragraph{What is and is not certified.}
The generators must generate \(G\), or the \emph{joint} image of
\((\rho,\sigma)\). Generating the image of \(\rho\) alone does not
suffice: an element in its kernel might act nontrivially under \(\sigma\).
For continuous representations it is enough to topologically generate
the joint image, since the equations are continuous and congruence classes
are closed. An arbitrary list of Frobenius elements is not automatically
such a generating set. The word bound above is an elementary finite
certificate, not a claim of an optimal matrix-algebra length bound.

\subsection{Frobenius traces give upper bounds without scalar elements}

For continuous Galois representations of a number field unramified
outside a finite set, Chebotarev density and continuity give
\begin{equation}
\label{eq:frobdepth}
 a=\min_{v\ {\rm unramified}}
 v_\pi\!\left(\operatorname{tr}\rho(\mathrm{Frob}_v)
             -\operatorname{tr}\sigma(\mathrm{Frob}_v)\right).
\end{equation}
Indeed, a finite minimal valuation defines a nonempty conjugacy-invariant
open subset in a finite quotient; some unramified Frobenius belongs to it.
Conversely every Frobenius is a group element. If all discrepancies are
zero, both sides are infinite.

One nonzero exact Frobenius discrepancy of valuation \(b\) proves
\(C\leq b\). Finding it needs no scalar Frobenius. To prove \(C=b\), one
must also establish congruence modulo \(\pi^b\), for example by an
invertible intertwiner, a proven joint-image certificate, or
Theorem~\ref{thm:finite}. Numerical approximations to Hecke eigenvalues are
not exact congruence certificates.

\subsection{The product over coefficient primes}
\label{sec:ideal}

The source's product-formula argument pays for local denominator losses
with logarithms of norms. Let \(F\) be a number field and \(S\) a finite
set of its finite places. Suppose a fixed compatible pair has trace
discrepancies \(\delta_v\in O_{F,S}\), with local representations integral
and of full matrix span at the coefficient places designated good,
all of which lie outside \(S\). Each chosen Frobenius place must be
unramified for every local pair under consideration; in particular,
coefficient places over its residue characteristic may have to be excluded.
These are additional arithmetic hypotheses, not consequences of
compatibility alone.

Let \(J=(\delta_{v_1},\ldots,\delta_{v_t})\) be a nonzero integral
\(O_{F,S}\)-ideal from exact discrepancies at finitely many common
unramified places. Write
\[
 N_S(J)=\prod_{\lambda\notin S}(N\lambda)^{v_\lambda(J)}.
\]
Then
\begin{equation}
\label{eq:idealbound}
 C_\lambda\leq v_\lambda(J),\qquad
 \sum_{\lambda\ {\rm good}}C_\lambda\log N\lambda
 \leq \log N_S(J).
\end{equation}
This follows from~\eqref{eq:frobdepth} and
\(v_\lambda(J)=\min_i v_\lambda(\delta_{v_i})\).
If \(J=O_{F,S}\), all these good-place global congruence terms vanish.
If the common discrepancies range over a Frobenius set that is
Chebotarev-dense for every good local pair, the ideal \(I\) of all
those discrepancies satisfies \(C_\lambda=v_\lambda(I)\) at good
places. This density is a separate requirement: simultaneous
unramifiedness can otherwise leave too small a common indexing set.
Noetherianity supplies a finite generating subfamily of \(I\); it does
not make an arbitrary sample complete.

This certifies the \emph{global congruence term}, not the whole Selmer
group. Local terms, bad coefficient primes, and the construction of the
compatible system remain separate. Discrepancy denominators must be
cleared and accounted for, or their primes included in \(S\).
Unrelated coefficient fields or unrelated twists cannot be combined
without first establishing a common field and a compatible family.

\paragraph{An explicit algebraic example.}
Over \(\mathbb Z[1/5]\), let \(G\) be free on \(u,v,z\), and put
\[
 A=\begin{pmatrix}1&1\\0&1\end{pmatrix},\qquad
 B=\begin{pmatrix}1&0\\1&1\end{pmatrix}.
\]
Set \(\rho(u)=\sigma(u)=A\), \(\rho(v)=\sigma(v)=B\),
\(\rho(z)=I\), and \(\sigma(z)=25I\).
The matrices \(I,A,B,AB\) form an integral basis, with trace-pairing
determinant \(-1\). Trace differences at \(z,zu,zv,zuv\) are
\(-48,-48,-48,-72\), generating \((24)\).
All discrepancies are divisible by \(24\) in \(\mathbb Z[1/5]\),
including inverse words. Thus \(C_2=3\), \(C_3=1\), and \(C_p=0\)
for \(p\notin\{2,3,5\}\). The exact denominator product is
\[
 2^3\,3=24,
\]
where the rank-two scalar estimate gives \(24^2=576\).
The relation matrix has Smith factors \(1,1,1,24\).
This is an algebraic model, not a claimed arithmetic Galois realization
or admissible \texttt{SelmerDatum}.
